% !PRINTED 591
\weberpage{591}
%<<CONT
come also the associated numbers. From 5. it then further follows
that a prime $p$ can be decomposed in only \gesperrt{one way} into the sum of two
squares. By this decomposition, however, the numbers $a, b$ are
not yet completely determined; rather, they can still be interchanged
with one another and furnished with two signs.
This comes down to saying that one can replace $\pi$ by
any of the four associated numbers $a+bi, -a-bi, -b+ai,$
$b-ai$. If $p$ is odd, then of the two
numbers $a, b$ one must be even, the other odd. If we choose,
say, for $b$ the even one of the two numbers, and determine
the sign so that $a \equiv 1 \pmod{4}$, then among the
four associated numbers a definite one is selected, which one may
\gesperrt{call the primary one}\footnote{Gauss gives at the place mentioned two different determinations
 of the primary numbers to choose from, of which this is the first.
He afterwards retains the second.}.

This definition of the primary numbers can be carried over to all
integers of the field $R(i)$ whose norm is
odd, and one then has the law that the product
of two primary numbers again yields a primary number.

The system of four associated numbers in the field $R(i)$ is
analogous to the pair of opposite numbers in the field $R$. In $R$
one regards the positive numbers as the primary ones.

By the formula (27) of the preceding section, for
any given prime $p$ one of the factors $a+bi$ is derived
from the division of the circle. By the formulae \S.\,172, (38), (39)
this number is primary if $p \equiv 5 \pmod{8}$, but opposite to
the primary one if $p \equiv 1 \pmod{8}$. But as to
which of the two conjugate numbers $a \pm bi$ is represented by these
formulae, we have no general criterion.

For example, the complex primes in $R(i)$ whose
norms are less than 200 are:
\[
\begin{aligned}
&1+i,\ 1+2i,\ 3+2i,\ 1+4i,\ 5+2i,\ 1+6i,\ 5+4i,\\
&7+2i,\ 5+6i,\ 3+8i,\ 5+8i,\ 9+4i,\ 1+10i,\ 3+10i,\\
&7+8i,\ 11+4i,\ 7+10i,\ 11+6i,\ 13+2i,\ 9+10i,\ 7+12i,\\
&1+14i.
\end{aligned}
\]
